
\begin{abstract}
We study population harvesting under a strong Allee effect,
where a population below a certain size may continue to
decline even without harvesting. Two harvesting methods
are considered: constant harvesting (Model A) and
proportional harvesting (Model B). We derive the equilibrium
points, analyse their stability, and calculate the critical
harvesting rates beyond which no positive equilibrium exists.
We also examine how the initial population affects survival,
calculate extinction times, and study delayed collapse near
the critical harvesting rate. The effect of changing
harvesting rates and restocking on population recovery is
also investigated. Numerical results are checked using
independent integration methods. We then apply
Physics-Informed Neural Networks (PINNs) to solve the
population equations and estimate unknown parameters from
synthetic observations. For the forward problem, the PINN
achieves relative $L^2$ errors ranging from
$2.9\times10^{-7}$ to $4.1\times10^{-3}$, although
conventional ODE solvers are considerably faster. For the
inverse problem, the PINN recovers all three fitted parameters
within $5\%$ of their true values in 10 out of 18
noise-free runs, compared with 18 out of 18 for ODE
least-squares fitting. We also study the effect of
measurement noise and limitations of parameter estimation,
including nonidentifiability in Model B when all four
physical parameters are unknown. All experiments use
synthetic data, and the results describe the mathematical
models rather than any particular real-world population.
\end{abstract}
